% =============================================================================
% Conclusion — workshop variant. Reframed per GPT-5.5 Pro v5 review (sec. C
% Move D): replaces vague open-question prose with 3-5 SHARP, CONCRETE
% workshop questions that reviewers/attendees could discuss.
% Target: <= 0.5 page.
% =============================================================================
\section{Discussion and Concrete Workshop Questions}
\label{sec:conclusion}

The workshop release is deliberately exploratory. Its strongest evidence is
short-horizon GSM8K dynamics; non-math coverage is a stress case, not a
result. We treat ZVF as a cheap descriptive diagnostic, not a predictive
failure detector, and we emphasize that the Qwen3-8B held-out control gives
a small, non-significant GRPO improvement over the base.

\paragraph{Five concrete questions for the workshop.}
We invite attendees to engage with five specific, testable questions raised
by this release rather than open-ended philosophizing:
\begin{enumerate}[leftmargin=1.4em,topsep=2pt,itemsep=2pt]
  \item Does step-level ZVF predict reward change \emph{within} a single
        run, controlling for run fixed effects, on stacks beyond Tinker?
        Our residualized within-GSM8K scatter
        (Figure~\ref{fig:workshop-killer}c) finds $r_{\mathrm{within}}$
        close to zero; is this a backend artifact or a general property?
  \item Under a fixed rollout-token budget, where does group size trade early
        step count for endgame contrast? Our two-seed $G{=}2\times160$ versus
        $G{=}16\times20$ panel ends with the small-$G$ arms at an all-correct
        ZVF wall while $G{=}16$ remains mid-learning; a separate 3-seed
        reanalysis places a fitted apex near $G{\approx}32$ for one stack and
        budget. Neither identifies a universal optimum. The needed experiment
        is a directly trained, multi-seed, token-matched sweep on open stacks.
  \item Are there tasks where high-ZVF cells \emph{do} predict downstream
        collapse, beyond GSM8K? Tool-use is a degenerate $\mathrm{ZVF}{=}1$
        case from step 1; the genuinely interesting test is on
        intermediate-density tasks (MATH-500 at low temperature, code
        with partial-credit rewards).
  \item What minimum disclosure makes a managed-API result count as
        benchmark evidence? Our Tinker rows are reproducible only at the
        level of API call scripts; should the community require GRPO
        loss / reward normalization / hardware metadata for a managed-API
        row to count as anything beyond a case study?
  \item Should ZVF-style diagnostic claims require raw per-group rollout
        JSONL by default? Our release bundles step-level aggregates for
        every run but raw JSONL only for a documented subset; we invite
        the workshop to align on a minimum raw-artifact bar.
\end{enumerate}
